\section{Introduction}
\label{sec:tpc46-introduction}

The TPC sequence isolates a deterministic equality-column problem whose
large-prime part is governed by
\begin{equation}
 pr-mj=h_0.
 \label{eq:tpc46-intro-determinant}
\end{equation}
Here \(p\) is prime, \(r\) is the complementary factor carrying a
M\"obius weight, \(m\) is an actual source row carrying a second
arithmetic coefficient, and \(j\) is an orbit variable with a smooth
profile restricted to a fixed residue class.  TPC--45 identifies the
absence of a third arbitrary
determinant-coefficient slot in the available fixed-determinant
estimates as the next analytic obstruction \citep{WangTPC45}.  The
natural response is to average the determinant shift and use the
progression-smooth \(j\)-slot.  The purpose of this paper is to
determine exactly
what that operation proves, which normalization it controls, and why
it does not yet return the physical column at \(h_0\).

There are three families that must not be merged.  A \emph{frozen
completed family} keeps one coefficient tensor and lets the
determinant value range through consecutive integers.  The inherited
TPC alias bank keeps one coefficient tensor but samples that family on
the progression

\[
 h_0+kM_{\rm al},
 \qquad
 M_{\rm al}=X^{399/400+o(1)}.
\]

A family reconstructed at different physical shifts changes the
eligible rows, residue classes, coprimality conditions, and generally
the coefficient tensor itself.  Fourier analysis applies directly to
the first family.  It applies to the second only after alias folding,
and it does not become a theorem about the third without a uniform
provenance statement.

\subsection{Main results}

Our first result is an exact positive closure theorem at a deliberately
named norm.  For scalar separated coefficients, put
\[
 \cT_h=\sum_{pr-mj=h}a_pb_rc_m\eta_j,
 \qquad
 \cD_\otimes=
 \left(\sum_{p,r}|a_pb_r|^2\right)
 \left(\sum_{m,j}|c_m\eta_j|^2\right).
\]
Multiplicative synthesis turns \(\cT_h\) into an additive correlation.
The resulting Cauchy--Schwarz estimate is
\begin{equation}
 |\cT_h|^2\le \rho_{pr}\rho_{mj}\cD_\otimes.
 \label{eq:tpc46-intro-tensor-bound}
\end{equation}
At the large-prime TPC interface, \(\rho_{pr}=1\) and
\(\rho_{mj}=X^{o(1)}\).  A finite Fourier--Bohr transform gives the
same tensor-scale closure after retaining the residual label
\(s=d(m)r\).  These are genuine fixed-shift theorems, but
\begin{equation}
 \cD_\otimes=\sum_h\Delta_h,
 \label{eq:tpc46-intro-diagonal-sum}
\end{equation}
where \(\Delta_h\) is the atomic diagonal at shift \(h\).  The
physical problem is normalized by \(\Delta_{h_0}\), not its sum over
all shifts.  Flat boxes and the prescribed-shift construction in
\cref{sec:tpc46-atomic-boundaries} show that the two scales can differ
by a power of \(X\).  Hence
\eqref{eq:tpc46-intro-tensor-bound} is not renamed as an
atomic-normalized TPC estimate.

Our second result is a deterministic shift-spectrum theorem.  For a
Hilbert-valued frozen packet with a \(j\)-independent vector coefficient,
the shift Fourier transform factors through
\begin{equation}
 K_m(\alpha)
 =\sum_{j\in\Z}w(j/J)\e(mj\alpha)
 =J\sum_{k\in\Z}\widehat w\bigl(J(k-m\alpha)\bigr).
 \label{eq:tpc46-intro-Poisson-kernel}
\end{equation}
If \(\operatorname{supp}\widehat w\subset[-\sigma,\sigma]\), this
kernel vanishes in the exact annulus
\begin{equation}
 \frac{\sigma}{M_-J}
 <\|\alpha\|_{\T}
 <\frac{1-\sigma/J}{M_+}.
 \label{eq:tpc46-intro-gap}
\end{equation}
At the endpoint this is
\[
 X^{-1+o(1)}<\|\alpha\|_{\T}<X^{-267/400+o(1)}.
\]
Restriction to a fixed progression in \(j\) changes only the constant
at the first alias and preserves these endpoint exponents.  Compactly
supported smooth orbit profiles satisfy rapid decay on every fixed
power-shrinking of the same annulus.  As consequences, suitably
supported band-pass shift averages vanish exactly, while sufficiently
coarse low-pass averages admit an exact zero-mode formula.  The smooth
version has a power-saving Poisson remainder.  We also compute the
corresponding averaged atomic diagonal and keep it distinct from the
full squared synthesis.

Our third result is a set of sharp transfer boundaries.  Weighted
averaging in \(h\) leaves the exact fiber-grouping norm equal to the
largest realized fiber when numerator and atomic diagonal use the same
weights.  Sampling only the TPC alias progression folds the entire
shift spectrum into \(M_{\rm al}\) branches.  For actual reconstructed
shifts, zero extension always gives a common ambient coordinate space,
while the present rigidity argument proves collision-free equality
labels throughout
\begin{equation}
 |h-h_0|\le
 X^{29629/210000-\epsilon},
 \label{eq:tpc46-intro-local-range}
\end{equation}
and makes no assertion that this sufficient range is optimal.  Resolving
the first Poisson alias by the present route needs a dense window of
length at least \(X^{267/400+o(1)}\).  Finally, positivity alone costs
the full effective number of shifts when an averaged nonnegative
quantity is localized back to one chosen shift.  These statements
separate four missing inputs: orbit-frequency regularity for the
pulled-back dual field, atomic rather than Cartesian normalization,
anti-alias control, and fixed-shift localization.

\subsection{Why existing trilinear estimates do not fill the gap}

Bettin--Chandee's trilinear Kloosterman estimate permits three arbitrary
sequences, but one is a Kloosterman-numerator sequence
\citep{BC2018}.  Their fixed-determinant corollary retains two smooth
determinant slots and only two arbitrary arithmetic slots.  Averaging
the determinant in the Bettin--Chandee--Radziwi\l\l application
compresses a frequency and a determinant into a structured numerator
sequence after smooth Poisson analysis; it does not remove those two
smooth-slot hypotheses \citep{BCR2017}.  The modulus-dispersion results
of Fouvry--Radziwi\l\l and Wright concern different averaged objects
\citep{FR2022,Wright2026}.  Moreover, the physical TPC dualization
introduces the joint field
\[
 z_{\text{output row},\,j,\,d(m)r},
\]
which need not split into their scalar coefficient slots.  The precise
comparison is given in \cref{sec:tpc46-literature-boundary}.

\subsection{Scope and organization}

All structural identities, finite operator norms, Fourier support
statements, Poisson formulas, and rational exponent comparisons stated
as theorems are proved without a probabilistic M\"obius hypothesis.
Statements about the actual TPC packet use only the inherited
coefficient interfaces explicitly cited from TPC--36, TPC--42,
TPC--43, and TPC--45 \citep{WangTPC36,WangTPC42,WangTPC43,WangTPC45}.
No theorem asserts that the arbitrary dual field is orbit-band-limited,
that the sparse alias bank can be unfolded at no cost, or that a shift
average controls the physical column.

\Cref{sec:tpc46-interface} reconstructs the exact TPC packet and its
dual coupling.  \Cref{sec:tpc46-model,sec:tpc46-tensor-closure,sec:tpc46-bohr}
prove the scalar and residual tensor closures.
\Cref{sec:tpc46-shift-spectrum,sec:tpc46-poisson-averages} establish the
spectral-gap and averaging formulas.
\Cref{sec:tpc46-atomic-boundaries,sec:tpc46-alias-localization} prove
the sharp normalization, folding, geometry, and localization
boundaries.  \Cref{sec:tpc46-literature-boundary} audits the available
external estimates, and \cref{sec:tpc46-next-gate} states four
reversible next doors.  The exact finite certificate and final claim
ledger appear in \cref{sec:tpc46-certificate}.

No result in this paper proves a fixed-shift TPC closure, a
Hardy--Littlewood prime-pair asymptotic, a breach of the sieve parity
barrier, or the twin-prime conjecture.
